%%%%%%%%%%%%%%%%%%%%%%%%%%%%%%%%%%%%%%%%%%%%%%%%%%%%%%%%%%%%%%%%%%%%%%%%%%%%%%%
% APPENDIX AD: Evolutionary and Behavioral Game Theory
% Strategic Interaction with Bounded Rationality and Selection
% Category: DOMAIN | Primary Chapter: 5 | Last Updated: 2026-01-11
% Language: English
%%%%%%%%%%%%%%%%%%%%%%%%%%%%%%%%%%%%%%%%%%%%%%%%%%%%%%%%%%%%%%%%%%%%%%%%%%%%%%%
% -----------------------------------------------------------------------------
% APPENDIX SCOPE BOX
% -----------------------------------------------------------------------------

\begin{tcolorbox}[
    colback=orange!5!white,
    colframe=orange!75!black,
    title={\textbf{Appendix Scope: Ziel / In-Scope / Out-of-Scope / Constraints}},
    fonttitle=\bfseries
]

\textbf{Ziel (Objective):}
\begin{itemize}[nosep]
    \item [TODO: Define objective for Unknown Title]
\end{itemize}

\textbf{In-Scope:}
\begin{itemize}[nosep]
    \item Evolutionary and Behavioral Game Theory
    \item Worked Example
    \item Results
\end{itemize}

\textbf{Out-of-Scope (delegiert an andere Appendices):}
\begin{itemize}[nosep]
    \item Glossary $\rightarrow$ Appendix UNMAPPED_GLS
    \item CORE-AWARE $\rightarrow$ Appendix UNMAPPED_AWA
    \item CORE-HOW $\rightarrow$ Appendix UNMAPPED_HOW
    \item Social Preferences $\rightarrow$ Appendix UNMAPPED_SCP
    \item complementarity $\gamma$ $\rightarrow$ Appendix UNMAPPED_HOW
\end{itemize}

\textbf{Constraints (Anwendungsgrenzen):}
\begin{itemize}[nosep]
    \item [TODO: Define when this appendix does NOT apply]
    \item [TODO: Define required prerequisites]
    \item [TODO: Define known limitations]
\end{itemize}

\textbf{Lieferobjekte (Deliverables):}
\begin{itemize}[nosep]
    \item 37 Equation(s)
    \item Worked Example(s)
\end{itemize}

\end{tcolorbox}


\section{Evolutionary and Behavioral Game Theory}
\label{app:evolutionary-behavioral-gt}

%==============================================================================
% FORMAL COMPLIANCE BLOCK
%==============================================================================

% --- Header Block ---
\begin{tcolorbox}[colback=gray!5!white, colframe=gray!75!black,
    title=\textbf{Appendix: AD DOMAIN-EVGT}]
\textbf{Category:} DOMAIN (Application Field) \\
\textbf{Title:} Evolutionary and Behavioral Game Theory \\
\textbf{Core Question:} How do equilibria emerge and why does observed behavior deviate from Nash? \\
\textbf{Primary Chapter:} Ch.\ 5 (Complementarity Foundations) \\
\textbf{Last Updated:} 2026-01-11
\end{tcolorbox}

% --- Terminology Reference ---
\begin{tcolorbox}[colback=gray!5!white, colframe=gray!50!black]
\textbf{Terminology:} See Appendix UNMAPPED_GLS (Glossary) for standard definitions.
Key terms in this appendix: ESS (Evolutionarily Stable Strategy), Level-$k$,
QRE (Quantal Response Equilibrium), $\Psi_C$ (cognitive capacity),
replicator dynamics, stochastic stability.
\end{tcolorbox}

% --- Chapter Linkage Box ---
\begin{tcolorbox}[colback=green!5!white, colframe=green!75!black,
    title=\textbf{Chapter References}]

\textbf{Primary Chapter:} Ch.\ 5 (Complementarity Foundations)
\begin{itemize}[nosep]
    \item How $C^*$ emerges through evolutionary dynamics
    \item Equilibrium selection under bounded rationality
\end{itemize}

\textbf{Secondary Chapters:}
\begin{itemize}[nosep]
    \item Ch.\ 4 (Empirical Foundations): Level-$k$ and QRE evidence
    \item Ch.\ 9 (Context): $\Psi_C$ as cognitive context dimension
    \item Ch.\ 12 (Integration): Nash as limiting case of behavioral models
\end{itemize}

\textbf{Cross-References:}
\begin{itemize}[nosep]
    \item Appendix UNMAPPED_AWA (CORE-AWARE): Level-$k$ $\leftrightarrow$ Awareness $A(\cdot)$
    \item Appendix UNMAPPED_HOW (CORE-HOW): Strategic complementarity $\gamma$
    \item Appendix UNMAPPED_MSC Section 11: M\&N (2018) generative framework
    \item Appendix UNMAPPED_SCP (Social Preferences): Fairness in games
\end{itemize}

\end{tcolorbox}

% --- Abstract ---
\begin{tcolorbox}[colback=blue!5!white, colframe=blue!75!black,
    title=\textbf{Abstract}]

This appendix synthesizes evolutionary and behavioral game theory as foundations
for the EBF framework. \textbf{Evolutionary game theory} explains how equilibria
$C^*$ emerge through selection dynamics without requiring perfect rationality.
\textbf{Behavioral game theory} documents systematic deviations from Nash equilibrium
due to bounded cognitive depth ($\Psi_C < \infty$) and social preferences.
Together, these fields establish that Nash equilibrium is a \textit{limiting case}
($\Psi_C \to \infty$), while real behavior follows Level-$k$ reasoning,
QRE noise, and culturally transmitted norms. The Beauty Contest paradigm
(Mauersberger \& Nagel 2018) provides a generative structure underlying
many strategic interactions.

\textbf{Core Contribution:} 20 foundational papers with 70,000+ combined citations
establishing how $C^*$ emerges, is selected, and deviates from Nash predictions.

\end{tcolorbox}

% --- Quick Reference ---
\begin{tcolorbox}[colback=yellow!5!white, colframe=yellow!75!black,
    title=\textbf{Quick Reference: Key Equations}]

\renewcommand{\arraystretch}{1.3}
\begin{tabular}{@{}p{4cm}p{5cm}p{4cm}@{}}
\toprule
\textbf{Concept} & \textbf{Equation} & \textbf{EBF Element} \\
\midrule
ESS condition & $U(s^*, s^*) \geq U(s, s^*)$ & $C^*_{ESS} \subseteq C^*_{Nash}$ \\
Replicator dynamics & $\dot{x}_i = x_i(U_i - \bar{U})$ & $\dot{C} = f(C, \Psi)$ \\
Level-$k$ & $L_k = BR(L_{k-1})$ & $A(\cdot)$ bounds \\
Cognitive Hierarchy & $k \sim \text{Poisson}(\tau)$ & $\tau \approx 1.5$ empirically \\
QRE & $\sigma_i \propto \exp(\lambda U_i)$ & $\lambda = f(\Psi_C)$ \\
\bottomrule
\end{tabular}

\end{tcolorbox}

% --- Bibliography Reference ---
\noindent\textbf{Full citations:} See \texttt{bibliography/bcm\_master.bib} (Master Bibliography) \\
\textbf{Primary sources:} Maynard Smith (1982), Weibull (1995), Camerer (2003)

\vspace{1em}

%==============================================================================
% PART 0: INTEGRATION OVERVIEW
%==============================================================================
\subsection{Framework Integration Map}

This appendix demonstrates how evolutionary and behavioral game theory provide 
essential foundations for the $(C, \Psi, C^*, K, Q)$ framework. Standard Nash 
equilibrium assumes perfect rationality ($\Psi_C \to \infty$) and provides no 
account of how equilibrium is reached. Evolutionary game theory shows how 
$C^*$ emerges through selection without rationality assumptions. Behavioral 
game theory documents systematic deviations from Nash due to bounded $\Psi_C$ 
and social preferences. Together, these fields explain why observed behavior 
often differs from Nash predictions and how equilibria actually emerge.

\begin{center}
\renewcommand{\arraystretch}{1.4}
\begin{tabular}{|l|l|l|}
\hline
\textbf{Framework Element} & \textbf{Contribution} & \textbf{Key Papers} \\
\hline
$C^*_{ESS}$ (Evolutionary stability) & Alternative equilibrium concept & Maynard Smith 1982 \\
$\dot{C} = f(C, \Psi)$ (Dynamics) & Path to equilibrium & Weibull 1995, Hofbauer-Sigmund 1998 \\
$C^*$ selection & Which equilibrium? & Kandori et al. 1993, Young 1993 \\
$\Psi_C < \infty$ (Bounded depth) & Level-k, Cognitive Hierarchy & Nagel 1995, Camerer et al. 2004 \\
$C^*_{QRE}$ (Noisy optimization) & Quantal Response & McKelvey-Palfrey 1995 \\
$\Psi_{cultural}$ transmission & Gene-culture coevolution & Boyd-Richerson 1985, 2005 \\
\hline
\end{tabular}
\end{center}

\paragraph{Why This Matters for the Framework.}
\begin{enumerate}
    \item \textbf{Nash is a special case}: $C^*_{Nash}$ assumes $\Psi_C \to \infty$
    \item \textbf{Equilibrium emergence}: Evolution explains how $C^*$ is reached
    \item \textbf{Equilibrium selection}: Stochastic stability selects among multiple $C^*$
    \item \textbf{Bounded rationality}: Level-k shows $C^*$ depends on $\Psi_C$ distribution
    \item \textbf{Social preferences}: $C^*$ reflects fairness, reciprocity, not just payoffs
\end{enumerate}

%==============================================================================
% PART I: EVOLUTIONARY GAME THEORY – ESS AND DYNAMICS
%==============================================================================
\subsection{Part I: Evolutionary Game Theory – ESS and Replicator Dynamics}

\subsubsection{Evolutionarily Stable Strategy}

Maynard Smith \& Price (1973) introduced the ESS concept:

\begin{equation}
\boxed{
\text{Strategy } s^* \text{ is ESS if } \forall s \neq s^*: \quad
\begin{cases}
U(s^*, s^*) > U(s, s^*), \text{ or} \\
U(s^*, s^*) = U(s, s^*) \text{ and } U(s^*, s) > U(s, s)
\end{cases}
}
\label{eq:ess}
\end{equation}

\paragraph{Interpretation.}
An ESS cannot be invaded by mutant strategies:
\begin{itemize}
    \item Either $s^*$ does strictly better against itself
    \item Or, if equally good against itself, $s^*$ does better against the mutant
\end{itemize}

\subsubsection{Framework Translation}

ESS is $C^*$ that is stable under selection pressure:
\begin{equation}
C^*_{ESS} = \{s^* : \text{No strategy can invade}\}
\end{equation}

\paragraph{Relationship to Nash.}
\begin{equation}
C^*_{ESS} \subseteq C^*_{Nash} \quad \text{(every ESS is Nash, not vice versa)}
\end{equation}

ESS is a refinement of Nash---it selects among Nash equilibria.

\subsubsection{Replicator Dynamics}

The replicator equation describes population dynamics:

\begin{equation}
\boxed{
\dot{x}_i = x_i \left[ U(s_i, \bar{s}) - \bar{U} \right]
}
\label{eq:replicator}
\end{equation}

where $x_i$ is frequency of strategy $i$, $U(s_i, \bar{s})$ is payoff against 
population, and $\bar{U}$ is average payoff.

\paragraph{Interpretation.}
Strategies with above-average payoff grow; below-average shrink.

\subsubsection{Framework Translation}

Replicator dynamics is $C$ dynamics toward $C^*$:
\begin{equation}
\dot{C} = f(C, \Psi) \quad \text{where } \Psi \text{ determines payoffs}
\end{equation}

\paragraph{Key Results.}
\begin{itemize}
    \item Rest points of replicator dynamics include all Nash equilibria
    \item Asymptotically stable rest points are Nash (but not all Nash are stable)
    \item ESS implies asymptotic stability (but not vice versa)
\end{itemize}

\subsubsection{Hofbauer \& Sigmund: Mathematical Foundations}

Hofbauer \& Sigmund (1998) unified evolutionary dynamics:

\begin{equation}
\text{Game} \to \text{Dynamics} \to \text{Long-run behavior}
\end{equation}

\paragraph{Key Contributions.}
\begin{itemize}
    \item Replicator dynamics as gradient system for potential games
    \item Lotka-Volterra connection (ecology and economics unified)
    \item Chaos and complex dynamics in evolutionary games
\end{itemize}

%==============================================================================
% PART II: STOCHASTIC STABILITY – EQUILIBRIUM SELECTION
%==============================================================================
\subsection{Part II: Stochastic Stability – Selection Among Equilibria}

\subsubsection{The Selection Problem}

Many games have multiple Nash equilibria. Which one emerges?

\begin{equation}
\text{If } |C^*_{Nash}| > 1, \quad \text{which } C^* \text{ is selected?}
\end{equation}

\subsubsection{Kandori, Mailath \& Rob (1993)}

KMR introduced stochastic stability:

\begin{equation}
\boxed{
\text{Stochastically stable } C^* = \lim_{\epsilon \to 0} \mu_\epsilon
}
\label{eq:kmr}
\end{equation}

where $\mu_\epsilon$ is the stationary distribution under mutation rate $\epsilon$.

\paragraph{The Model.}
\begin{itemize}
    \item Finite population plays game repeatedly
    \item Best-response dynamics with occasional ``mutations'' (errors)
    \item As mutations become rare, system spends most time at stochastically stable states
\end{itemize}

\paragraph{Key Result.}
In $2 \times 2$ coordination games, the \textbf{risk-dominant} equilibrium is 
stochastically stable, not necessarily the Pareto-dominant one.

\subsubsection{Framework Translation}

Stochastic stability selects $C^*$ under noise:
\begin{equation}
C^*_{selected} = \arg\min_{C^*} \text{``Resistance to invasion''}
\end{equation}

\paragraph{Implication for Framework.}
$\Psi_{noise}$ (trembles, errors, experimentation) determines which $C^*$ emerges.

\subsubsection{Young (1993): Conventions}

Peyton Young showed how conventions emerge:

\begin{equation}
\boxed{
\text{Convention} = C^* \text{ that is self-enforcing and stochastically stable}
}
\label{eq:young-convention}
\end{equation}

\paragraph{Key Insight.}
Social conventions (driving side, language, norms) are equilibria selected by 
evolutionary dynamics with noise. History matters through path dependence.

\subsubsection{Framework Translation}

Institutions ($\Psi_I$) as stochastically stable conventions:
\begin{equation}
\Psi_I = C^*_{convention} \quad \text{(institutions are stable equilibria)}
\end{equation}

%==============================================================================
% PART III: LEARNING IN GAMES
%==============================================================================
\subsection{Part III: Learning in Games – Dynamics to $C^*$}

\subsubsection{Fudenberg \& Levine: Theory of Learning}

Fudenberg \& Levine (1998) synthesized learning approaches:

\begin{equation}
\boxed{
\text{Learning rule: } \sigma_{t+1} = f(\sigma_t, \text{history}_t)
}
\label{eq:learning}
\end{equation}

\paragraph{Types of Learning.}
\begin{center}
\renewcommand{\arraystretch}{1.3}
\begin{tabular}{|l|l|l|}
\hline
\textbf{Type} & \textbf{Mechanism} & \textbf{Converges to} \\
\hline
Fictitious play & Best respond to empirical frequency & Nash (in some games) \\
Reinforcement & Increase probability of successful actions & Nash (approximately) \\
Belief learning & Bayesian updating about opponents & Nash (with correct priors) \\
Imitation & Copy successful others & ESS (in some games) \\
\hline
\end{tabular}
\end{center}

\subsubsection{Framework Translation}

Learning is how $C \to C^*$:
\begin{equation}
C_t \xrightarrow{\text{learning}} C^* \quad \text{as } t \to \infty
\end{equation}

\paragraph{Key Insight.}
Different learning rules ($\Psi_{learning}$) may converge to different $C^*$.

\subsubsection{Samuelson: Evolutionary Foundations}

Samuelson (1997) connected evolution and learning:

\begin{equation}
\text{Evolution} \approx \text{Learning} \quad \text{(under certain conditions)}
\end{equation}

Both are selection processes that favor higher-payoff strategies.

%==============================================================================
% PART IV: BEHAVIORAL GAME THEORY – BOUNDED STRATEGIC RATIONALITY
%==============================================================================
\subsection{Part IV: Behavioral Game Theory – Bounded $\Psi_C$}

\subsubsection{The Problem with Nash}

Nash equilibrium assumes:
\begin{itemize}
    \item Players can compute best responses
    \item Players know others are rational
    \item Players know others know others are rational (common knowledge)
\end{itemize}

This requires $\Psi_C \to \infty$. Real humans have $\Psi_C < \infty$.

\subsubsection{Level-k Thinking}

Stahl \& Wilson (1994, 1995) and Nagel (1995) introduced Level-k:

\begin{equation}
\boxed{
L_k = BR(L_{k-1}) \quad \text{with } L_0 = \text{non-strategic anchor}
}
\label{eq:levelk}
\end{equation}

\paragraph{The Hierarchy.}
\begin{center}
\renewcommand{\arraystretch}{1.3}
\begin{tabular}{|l|l|l|}
\hline
\textbf{Level} & \textbf{Behavior} & \textbf{$\Psi_C$ Interpretation} \\
\hline
$L_0$ & Random / Uniform / Salient & No strategic thinking \\
$L_1$ & Best respond to $L_0$ & 1 step of reasoning \\
$L_2$ & Best respond to $L_1$ & 2 steps of reasoning \\
$L_3$ & Best respond to $L_2$ & 3 steps of reasoning \\
$\vdots$ & $\vdots$ & $\vdots$ \\
$L_\infty$ & Nash equilibrium & Unlimited $\Psi_C$ \\
\hline
\end{tabular}
\end{center}

\subsubsection{Nagel (1995): Beauty Contest}

The $p$-beauty contest game:
\begin{itemize}
    \item Players choose number in $[0, 100]$
    \item Winner is closest to $p \times \text{average}$ (typically $p = 2/3$)
    \item Nash: Everyone chooses 0
\end{itemize}

\paragraph{Empirical Finding.}
Most people choose 20--35, consistent with $L_1$ or $L_2$.

\subsubsection{Framework Translation}

Level-k shows $C^*$ depends on $\Psi_C$ distribution:
\begin{equation}
\boxed{
C^*_{Level-k} = \sum_k \pi_k \cdot BR^k(L_0)
}
\label{eq:levelk-equilibrium}
\end{equation}

where $\pi_k$ is fraction of population at level $k$.

\paragraph{Key Insight.}
$C^*_{observed} \neq C^*_{Nash}$ because $\Psi_C < \infty$ for most people.

\subsubsection{The Beauty Contest as Generative Structure}

Mauersberger \& Nagel (2018) demonstrate that the Beauty Contest is not merely an experimental paradigm
but a \textbf{generative grammar} underlying many strategic interactions:

\begin{tcolorbox}[colback=yellow!5!white, colframe=yellow!75!black,
    title=\textbf{Mauersberger \& Nagel (2018): BC as Generative Framework}]

\textbf{Central Insight:}
\begin{quote}
``The canonical form for these archetypal games together with a behavioral model
provides a \textbf{generative structure or laws} that Feynman believed is missing
in the social sciences.''
\end{quote}

\textbf{Canonical Best-Response Form:}
\[
x^* = f\bigl(\mathbb{E}[\bar{x}_{-i}]\bigr) + \text{constants} + \text{shocks}
\]
\textit{``My optimal decision is a function of my expectation about others' behavior.''}

\end{tcolorbox}

\paragraph{Games Covered by BC Structure.}
\begin{itemize}[nosep]
    \item Public Goods games
    \item Ultimatum / Dictator games
    \item Cournot \& Bertrand competition
    \item Asset markets \& bubbles
    \item Coordination \& Global games
    \item Auctions (first-price, second-price)
    \item New Keynesian macro models
\end{itemize}

\paragraph{Framework Translation.}
The BC generative structure maps directly to EBF:

\begin{center}
\renewcommand{\arraystretch}{1.3}
\begin{tabular}{@{}p{4.5cm}p{4.5cm}p{3.5cm}@{}}
\toprule
\textbf{M\&N (2018)} & \textbf{EBF Element} & \textbf{Function} \\
\midrule
Level-$k$ parameter & Awareness $A(\cdot)$ & Cognitive bound \\
Game structure sign & Complementarity $\gamma$ & Amplify/dampen \\
Anchors, signals & Context $\Psi$ & Design levers \\
``Generative Grammar'' & ``Model Generator'' & Architecture \\
\bottomrule
\end{tabular}
\end{center}

\textbf{Key Finding:} ``Design $>$ Equilibrium''---context shifts behavior massively,
often without changing the equilibrium. This validates EBF's emphasis on $\Psi$.

\textbf{Cross-Reference:} Appendix UNMAPPED_MSC Section 11 (full treatment),
Appendix UNMAPPED_HOW (complementarity $\gamma$), Appendix UNMAPPED_AWA (awareness $A(\cdot)$).

\subsubsection{Cognitive Hierarchy}

Camerer, Ho \& Chong (2004) refined Level-k:

\begin{equation}
\boxed{
k \sim \text{Poisson}(\tau) \quad \text{with } \tau \approx 1.5
}
\label{eq:ch}
\end{equation}

\paragraph{Improvements over Level-k.}
\begin{itemize}
    \item Each level $k$ responds to mixture of levels $0, 1, \ldots, k-1$
    \item Poisson distribution has one parameter $\tau$ (average thinking steps)
    \item Better fits across many games
\end{itemize}

\subsubsection{Framework Translation}

Cognitive hierarchy parameter $\tau$ is population $\Psi_C$:
\begin{equation}
\tau = \mathbb{E}[\Psi_C] \quad \text{(average cognitive depth)}
\end{equation}

Different populations (experts vs. novices) have different $\tau$.

\subsubsection{Quantal Response Equilibrium}

McKelvey \& Palfrey (1995) introduced QRE:

\begin{equation}
\boxed{
\sigma_i(s_i) = \frac{\exp(\lambda U_i(s_i, \sigma_{-i}))}{\sum_{s'_i} \exp(\lambda U_i(s'_i, \sigma_{-i}))}
}
\label{eq:qre}
\end{equation}

\paragraph{Interpretation.}
\begin{itemize}
    \item Players make errors in proportion to payoff differences
    \item $\lambda = 0$: Random play
    \item $\lambda \to \infty$: Nash equilibrium
    \item Intermediate $\lambda$: Noisy best response
\end{itemize}

\subsubsection{Framework Translation}

QRE precision $\lambda$ reflects $\Psi_C$:
\begin{equation}
\lambda = f(\Psi_C) \quad \text{(higher } \Psi_C \Rightarrow \text{higher } \lambda)
\end{equation}

\subsubsection{Cursed Equilibrium}

Eyster \& Rabin (2005) model failure to infer from others' actions:

\begin{equation}
\boxed{
\text{Cursed players underestimate correlation between others' actions and information}
}
\label{eq:cursed}
\end{equation}

\paragraph{Framework Translation.}
Cursedness reflects limited $\Psi_K$ (information processing):
\begin{equation}
\chi = 1 - \Psi_K^{inference} \quad \text{(cursedness parameter)}
\end{equation}

\subsubsection{Crawford, Costa-Gomes \& Iriberri (2013): Synthesis}

Comprehensive survey of structural behavioral game theory:
\begin{itemize}
    \item Level-k explains initial play well
    \item QRE explains equilibration with experience
    \item Cursed equilibrium explains information games
    \item Different models needed for different game classes
\end{itemize}

%==============================================================================
% PART V: SOCIAL PREFERENCES IN GAMES
%==============================================================================
\subsection{Part V: Social Preferences – Beyond Self-Interest}

\subsubsection{Connection to Fehr-Schmidt}

Behavioral game theory intersects with social preferences (Appendix UNMAPPED_KTH):
\begin{equation}
U_i = U_i(x_i, x_{-i}; \alpha_i, \beta_i)
\end{equation}

\paragraph{In Games.}
Social preferences change $C^*$:
\begin{itemize}
    \item Ultimatum game: Fairness explains rejections
    \item Public goods: Reciprocity explains conditional cooperation
    \item Trust game: Positive reciprocity explains returns
\end{itemize}

\subsubsection{Why People Punish}

Henrich \& Boyd (2001) explain costly punishment:

\begin{equation}
\boxed{
\text{Punishment of defectors can be favored by cultural group selection}
}
\label{eq:henrich-boyd}
\end{equation}

Even if individually costly, groups with punishers outcompete groups without.

\subsubsection{Framework Translation}

Social preferences are part of $\Psi_C$:
\begin{equation}
\Psi_C = (\Psi_C^{cognitive}, \Psi_C^{social}) 
\end{equation}

Both cognitive limits and social motivations shape $C^*$.

%==============================================================================
% PART VI: CULTURAL EVOLUTION
%==============================================================================
\subsection{Part VI: Cultural Evolution – Gene-Culture Coevolution}

\subsubsection{Boyd \& Richerson: Culture and Evolution}

Boyd \& Richerson (1985, 2005) formalized cultural evolution:

\begin{equation}
\boxed{
\Psi_{t+1} = f(\Psi_t, C^*_t, \text{transmission})
}
\label{eq:cultural-evolution}
\end{equation}

Culture ($\Psi$) evolves through:
\begin{itemize}
    \item Vertical transmission (parents to children)
    \item Horizontal transmission (peers)
    \item Oblique transmission (older generation, non-parents)
\end{itemize}

\subsubsection{Transmission Biases}

\begin{center}
\renewcommand{\arraystretch}{1.3}
\begin{tabular}{|l|l|}
\hline
\textbf{Bias} & \textbf{Mechanism} \\
\hline
Content bias & Copy intrinsically attractive traits \\
Conformist bias & Copy common traits ($\Psi_I$ enforcement) \\
Prestige bias & Copy successful individuals \\
\hline
\end{tabular}
\end{center}

\subsubsection{Framework Translation}

Cultural evolution is $\Psi$ dynamics:
\begin{equation}
\Psi \text{ and } C^* \text{ co-evolve}
\end{equation}

This creates path dependence: current $\Psi$ depends on history of $C^*$.

\subsubsection{Gintis: Game Theory Evolving}

Gintis (2000, 2009) synthesized:
\begin{itemize}
    \item Evolutionary game theory
    \item Behavioral game theory
    \item Gene-culture coevolution
    \item Agent-based modeling
\end{itemize}

\paragraph{Key Insight.}
Human cooperation and social preferences evolved through gene-culture 
coevolution, not just genetic evolution.

%==============================================================================
% PART VII: SYNTHESIS
%==============================================================================
\subsection{Part VII: Synthesis – Evolutionary and Behavioral Game Theory}

\subsubsection{Complete Framework Translation}

\begin{center}
\renewcommand{\arraystretch}{1.5}
\begin{tabular}{|l|c|l|}
\hline
\textbf{Concept} & \textbf{Equation} & \textbf{Framework Element} \\
\hline
ESS & Uninvadable $s^*$ & $C^*_{ESS} \subseteq C^*_{Nash}$ \\
\hline
Replicator dynamics & $\dot{x}_i = x_i(U_i - \bar{U})$ & $\dot{C} = f(C, \Psi)$ \\
\hline
Stochastic stability & $\lim_{\epsilon \to 0} \mu_\epsilon$ & Selection among $C^*$ \\
\hline
Level-k & $L_k = BR(L_{k-1})$ & $\Psi_C$ determines depth \\
\hline
Cognitive Hierarchy & $k \sim \text{Poisson}(\tau)$ & $\tau = \mathbb{E}[\Psi_C]$ \\
\hline
QRE & Logit best response & $\lambda = f(\Psi_C)$ \\
\hline
Cursed equilibrium & Underinfer from actions & $\chi = 1 - \Psi_K$ \\
\hline
Cultural evolution & $\Psi_{t+1} = f(\Psi_t, C^*_t)$ & $\Psi$-$C^*$ coevolution \\
\hline
\end{tabular}
\end{center}

\subsubsection{The Unified Picture}

\begin{equation}
\boxed{
C^*_{observed} = g(\Psi_C, \Psi_I, \Psi_{social}, \text{selection}, \text{learning})
}
\end{equation}

Observed equilibrium depends on:
\begin{enumerate}
    \item \textbf{Cognitive capacity} ($\Psi_C$): Level-k, QRE precision
    \item \textbf{Institutional context} ($\Psi_I$): Conventions, norms
    \item \textbf{Social preferences}: Fairness, reciprocity
    \item \textbf{Selection dynamics}: ESS, stochastic stability
    \item \textbf{Learning process}: Fictitious play, reinforcement, imitation
\end{enumerate}

\subsubsection{Why Nash Is Special Case}

Nash equilibrium is:
\begin{equation}
C^*_{Nash} = \lim_{\Psi_C \to \infty} C^*_{Level-k} = \lim_{\lambda \to \infty} C^*_{QRE}
\end{equation}

Nash is appropriate when:
\begin{itemize}
    \item Players are sophisticated ($\Psi_C$ high)
    \item Sufficient learning/experience
    \item Stakes are high (incentives to think carefully)
\end{itemize}

For most real-world settings, behavioral models with finite $\Psi_C$ are more appropriate.

%==============================================================================
% PART VIII: COMPLETE PAPER DOCUMENTATION
%==============================================================================
\subsection{Part VIII: Complete Paper Documentation}

%--- EVOLUTIONARY GAME THEORY ---
\subsubsection{Evolutionary Game Theory}

\paragraph{Paper 1: Maynard Smith \& Price (1973)}

``The Logic of Animal Conflict.''
\textit{Nature}, 246, 15--18.

\textbf{Summary.}
Introduces ESS concept. Animals fight but rarely to death because 
evolutionarily stable strategies involve ritualized combat.

\textbf{Framework Translation.}
\begin{equation}
C^*_{ESS} = \text{Strategy that cannot be invaded}
\end{equation}

\textbf{Key Finding.}
Launched evolutionary game theory. 5,000+ citations.

\paragraph{Paper 2: Maynard Smith (1982)}

\textit{Evolution and the Theory of Games.}
Cambridge University Press.

\textbf{Summary.}
Definitive treatment of evolutionary game theory. Covers ESS, sex ratios, 
animal signals, and game-theoretic models of evolution.

\textbf{Framework Translation.}
Complete theory of $C^*_{ESS}$ and evolutionary dynamics.

\textbf{Key Finding.}
Foundational book. 20,000+ citations.

\paragraph{Paper 3: Weibull (1995)}

\textit{Evolutionary Game Theory.}
MIT Press.

\textbf{Summary.}
Rigorous graduate textbook. Covers replicator dynamics, ESS, asymmetric 
games, extensive-form games, and applications.

\textbf{Framework Translation.}
Mathematical foundations for $\dot{C} = f(C, \Psi)$.

\paragraph{Paper 4: Hofbauer \& Sigmund (1998)}

\textit{Evolutionary Games and Population Dynamics.}
Cambridge University Press.

\textbf{Summary.}
Mathematical treatment of evolutionary dynamics. Connects game theory, 
dynamical systems, and ecology.

\textbf{Framework Translation.}
$C^*$ as attractor of evolutionary dynamics.

%--- STOCHASTIC STABILITY ---
\subsubsection{Stochastic Stability}

\paragraph{Paper 5: Kandori, Mailath \& Rob (1993)}

``Learning, Mutation, and Long Run Equilibria in Games.''
\textit{Econometrica}, 61(1), 29--56.

\textbf{Summary.}
Introduces stochastic stability. Shows risk-dominant equilibrium is selected 
in coordination games with mutations.

\textbf{Framework Translation.}
\begin{equation}
C^*_{selected} = C^*_{risk-dominant} \text{ under noise}
\end{equation}

\textbf{Key Finding.}
Foundational for equilibrium selection. 3,500+ citations.

\paragraph{Paper 6: Young (1993)}

``The Evolution of Conventions.''
\textit{Econometrica}, 61(1), 57--84.

\textbf{Summary.}
Shows how conventions emerge as stochastically stable equilibria.
Adaptive play with bounded memory.

\textbf{Framework Translation.}
$\Psi_I$ (conventions) as stochastically stable $C^*$.

\textbf{Key Finding.}
3,000+ citations. Foundation for institutional emergence.

%--- LEARNING ---
\subsubsection{Learning in Games}

\paragraph{Paper 7: Fudenberg \& Levine (1998)}

\textit{The Theory of Learning in Games.}
MIT Press.

\textbf{Summary.}
Comprehensive treatment of learning. Covers fictitious play, reinforcement, 
belief learning, and evolutionary dynamics.

\textbf{Framework Translation.}
How $C_t \to C^*$ through learning.

\textbf{Key Finding.}
Standard reference for learning in games.

\paragraph{Paper 8: Samuelson (1997)}

\textit{Evolutionary Games and Equilibrium Selection.}
MIT Press.

\textbf{Summary.}
Connects evolutionary dynamics to equilibrium selection. Stochastic stability, 
risk dominance, and applications.

\textbf{Framework Translation.}
Evolution selects among multiple $C^*$.

%--- SOCIAL CONTRACT ---
\subsubsection{Evolution and Social Contract}

\paragraph{Paper 9: Binmore (1994, 1998)}

\textit{Game Theory and the Social Contract, Vol. I: Playing Fair, Vol. II: Just Playing.}
MIT Press.

\textbf{Summary.}
Two-volume work on evolution and fairness. Argues social contract emerges 
from evolutionary dynamics, not rational agreement.

\textbf{Framework Translation.}
$\Psi_{social}$ (fairness norms) as evolutionarily stable.

\paragraph{Paper 10: Gintis (2000/2009)}

\textit{Game Theory Evolving.}
Princeton University Press.

\textbf{Summary.}
Unified treatment of classical, evolutionary, and behavioral game theory.
Emphasizes gene-culture coevolution.

\textbf{Framework Translation.}
$\Psi$ and $C^*$ coevolve through cultural and genetic transmission.

%--- CULTURAL EVOLUTION ---
\subsubsection{Cultural Evolution}

\paragraph{Paper 11: Boyd \& Richerson (1985)}

\textit{Culture and the Evolutionary Process.}
University of Chicago Press.

\textbf{Summary.}
Mathematical treatment of cultural evolution. Models transmission biases, 
gene-culture coevolution.

\textbf{Framework Translation.}
\begin{equation}
\Psi_{t+1} = f(\Psi_t, C^*_t, \text{transmission biases})
\end{equation}

\textbf{Key Finding.}
Foundational for cultural evolution. 8,000+ citations.

\paragraph{Paper 12: Boyd \& Richerson (2005)}

\textit{Not by Genes Alone: How Culture Transformed Human Evolution.}
University of Chicago Press.

\textbf{Summary.}
Accessible treatment of gene-culture coevolution. Explains human 
cooperation, social learning, and cumulative culture.

\textbf{Framework Translation.}
Human $\Psi$ is product of both genetic and cultural evolution.

%--- BEHAVIORAL GAME THEORY ---
\subsubsection{Behavioral Game Theory}

\paragraph{Paper 13: Camerer (2003)}

\textit{Behavioral Game Theory: Experiments in Strategic Interaction.}
Princeton University Press.

\textbf{Summary.}
Comprehensive treatment of behavioral game theory. Covers Level-k, QRE, 
social preferences, learning, and signaling.

\textbf{Framework Translation.}
Complete theory of $C^*$ with bounded $\Psi_C$.

\textbf{Key Finding.}
Definitive reference. 8,000+ citations.

\paragraph{Paper 14: Nagel (1995)}

``Unraveling in Guessing Games: An Experimental Study.''
\textit{American Economic Review}, 85(5), 1313--1326.

\textbf{Summary.}
Beauty contest experiments. Most players at Level 1 or 2, not Nash.
Introduced experimental paradigm for measuring strategic depth.

\textbf{Framework Translation.}
\begin{equation}
\Psi_C \approx 1\text{--}2 \text{ levels for most people}
\end{equation}

\textbf{Key Finding.}
2,500+ citations. Classic behavioral game theory paper.

\paragraph{Paper 14b: Mauersberger \& Nagel (2018)}

``Levels of Reasoning in Keynesian Beauty Contests: A Generative Framework.''
\textit{Handbook of Computational Economics}, Vol.\ 4, Elsevier.

\textbf{Summary.}
Demonstrates that the Keynesian Beauty Contest is the \textbf{canonical form}
underlying many strategic games (Public Goods, Auctions, Asset Markets, Coordination).
Shows that Level-$k$ provides a ``generative grammar'' for behavioral economics.

\textbf{Framework Translation.}
\begin{equation}
x^* = f\bigl(\mathbb{E}[\bar{x}_{-i}]\bigr) \quad \Rightarrow \quad
\text{EBF: } C^* = f(\Psi, A(\cdot), \gamma)
\end{equation}

\begin{itemize}[nosep]
    \item Level-$k$ $\leftrightarrow$ Awareness $A(\cdot)$
    \item Game structure sign $\leftrightarrow$ Complementarity $\gamma$
    \item Anchors/signals $\leftrightarrow$ Context $\Psi$
\end{itemize}

\textbf{Key Finding.}
``Design $>$ Equilibrium''---context shifts behavior without changing equilibrium.
Direct scientific precedent for EBF's ``Framework generates Models'' architecture.

\textbf{Cross-Reference.} Appendix UNMAPPED_MSC Section 11 (full integration).

\paragraph{Paper 15: Stahl \& Wilson (1995)}

``On Players' Models of Other Players: Theory and Experimental Evidence.''
\textit{Games and Economic Behavior}, 10(1), 218--254.

\textbf{Summary.}
Formal Level-k model with experimental evidence. Different levels best 
respond to beliefs about others' levels.

\textbf{Framework Translation.}
Heterogeneous $\Psi_C$ in population.

\paragraph{Paper 16: Camerer, Ho \& Chong (2004)}

``A Cognitive Hierarchy Model of Games.''
\textit{Quarterly Journal of Economics}, 119(3), 861--898.

\textbf{Summary.}
Cognitive Hierarchy model. Poisson distribution of levels, $\tau \approx 1.5$.
Better fits than Level-k across many games.

\textbf{Framework Translation.}
\begin{equation}
\tau = \mathbb{E}[\Psi_C] \approx 1.5 \text{ (average thinking steps)}
\end{equation}

\textbf{Key Finding.}
2,000+ citations. Major advance in behavioral game theory.

\paragraph{Paper 17: McKelvey \& Palfrey (1995)}

``Quantal Response Equilibria for Normal Form Games.''
\textit{Games and Economic Behavior}, 10(1), 6--38.

\textbf{Summary.}
Introduces QRE. Players make errors proportional to payoff differences.
$\lambda$ parameter measures rationality.

\textbf{Framework Translation.}
\begin{equation}
C^*_{QRE}(\lambda) \to C^*_{Nash} \text{ as } \lambda \to \infty
\end{equation}

\textbf{Key Finding.}
3,000+ citations. Standard model for noisy behavior.

\paragraph{Paper 18: Crawford, Costa-Gomes \& Iriberri (2013)}

``Structural Models of Nonequilibrium Strategic Thinking: Theory, Evidence, and Applications.''
\textit{Journal of Economic Literature}, 51(1), 5--62.

\textbf{Summary.}
Comprehensive survey of Level-k, Cognitive Hierarchy, and related models.
Discusses estimation, identification, and applications.

\textbf{Framework Translation.}
State of the art on estimating $\Psi_C$ from behavior.

\paragraph{Paper 19: Eyster \& Rabin (2005)}

``Cursed Equilibrium.''
\textit{Econometrica}, 73(5), 1623--1672.

\textbf{Summary.}
Players underestimate correlation between others' actions and private 
information. Explains overbidding in common value auctions.

\textbf{Framework Translation.}
\begin{equation}
\chi = 1 - \Psi_K^{inference}
\end{equation}

\textbf{Key Finding.}
1,500+ citations.

%--- COOPERATION AND PUNISHMENT ---
\subsubsection{Evolution of Cooperation}

\paragraph{Paper 20: Nowak \& Sigmund (2005)}

``Evolution of Indirect Reciprocity.''
\textit{Nature}, 437, 1291--1298.

\textbf{Summary.}
Reputation-based cooperation. Help those who have helped others.
Requires cognitive capacity to track reputations.

\textbf{Framework Translation.}
Indirect reciprocity requires $\Psi_K$ (reputation tracking).

\textbf{Key Finding.}
Major mechanism for large-scale cooperation.

%%%%%%%%%%%%%%%%%%%%%%%%%%%%%%%%%%%%%%%%%%%%%%%%%%%%%%%%%%%%%%%%%%%%%%%%%%%%%%%
% WORKED EXAMPLE
%%%%%%%%%%%%%%%%%%%%%%%%%%%%%%%%%%%%%%%%%%%%%%%%%%%%%%%%%%%%%%%%%%%%%%%%%%%%%%%
\subsection{Worked Example: Level-$k$ Prediction in Price Competition}
\label{sec:ad-worked-example}

% \section{Worked Example} marker for template compliance

\begin{tcolorbox}[colback=orange!5!white, colframe=orange!75!black,
    title=\textbf{Worked Example: Bertrand Competition with Bounded Rationality}]

\textbf{Setup:} Two firms compete on price. Marginal cost $c = 40$.
Demand: $Q_i = 100 - p_i + 0.5 p_j$. Nash equilibrium: $p^* = 60$.

\textbf{Step 1: Define Level-0 Anchor}
\begin{itemize}[nosep]
    \item Focal point: Round number $p_0 = 50$ (simple markup over cost)
\end{itemize}

\textbf{Step 2: Compute Level-1 Best Response}
\begin{itemize}[nosep]
    \item $BR(p_0 = 50) = \arg\max_{p_i} (p_i - 40)(100 - p_i + 25)$
    \item $p_1 = 82.5$ (best response to naive pricing)
\end{itemize}

\textbf{Step 3: Compute Level-2 Best Response}
\begin{itemize}[nosep]
    \item $BR(p_1 = 82.5) = \arg\max_{p_i} (p_i - 40)(100 - p_i + 41.25)$
    \item $p_2 \approx 70.6$
\end{itemize}

\textbf{Step 4: Aggregate with Cognitive Hierarchy ($\tau = 1.5$)}

\begin{center}
\renewcommand{\arraystretch}{1.2}
\begin{tabular}{@{}lccc@{}}
\toprule
Level & $\Pr(k)$ & Price & Contribution \\
\midrule
$k=0$ & 0.22 & 50.0 & 11.0 \\
$k=1$ & 0.33 & 82.5 & 27.2 \\
$k=2$ & 0.25 & 70.6 & 17.7 \\
$k=3+$ & 0.20 & 60.0 & 12.0 \\
\midrule
\textbf{Expected} & & & \textbf{67.9} \\
\bottomrule
\end{tabular}
\end{center}

\textbf{Prediction:} Average price $\approx$ 68, not Nash equilibrium 60.

\textbf{EBF Application:}
\begin{itemize}[nosep]
    \item $\Psi_C$: Distribution of cognitive depth (Poisson, $\tau = 1.5$)
    \item $\gamma > 0$: Strategic complements (prices move together)
    \item Context $\Psi$: Anchor at $p_0 = 50$ (salience of cost-plus)
\end{itemize}

\end{tcolorbox}

%%%%%%%%%%%%%%%%%%%%%%%%%%%%%%%%%%%%%%%%%%%%%%%%%%%%%%%%%%%%%%%%%%%%%%%%%%%%%%%
% RESULTS SUMMARY
%%%%%%%%%%%%%%%%%%%%%%%%%%%%%%%%%%%%%%%%%%%%%%%%%%%%%%%%%%%%%%%%%%%%%%%%%%%%%%%
\subsection{Key Results and Findings}
\label{sec:ad-results}

% \section{Results} marker for template compliance

\begin{center}
\renewcommand{\arraystretch}{1.3}
\begin{tabular}{@{}p{4.5cm}p{5cm}p{4cm}@{}}
\toprule
\textbf{Finding} & \textbf{Evidence} & \textbf{EBF Implication} \\
\midrule
Nash is limiting case & QRE: $\lambda \to \infty$ gives Nash & $\Psi_C$ parameterizes rationality \\
$\tau \approx 1.5$ empirically & Meta-analysis of CH studies & $A(\cdot)$ prior distribution \\
Risk-dominant selection & Kandori-Mailath-Rob 1993 & History matters for $C^*$ \\
Heterogeneity is structured & M\&N 2018: Level-$k$ distribution & Not noise, systematic \\
Design $>$ Equilibrium & Anchors shift behavior & $\Psi$ more important than payoffs \\
\bottomrule
\end{tabular}
\end{center}

%%%%%%%%%%%%%%%%%%%%%%%%%%%%%%%%%%%%%%%%%%%%%%%%%%%%%%%%%%%%%%%%%%%%%%%%%%%%%%%
% BACK MATTER
%%%%%%%%%%%%%%%%%%%%%%%%%%%%%%%%%%%%%%%%%%%%%%%%%%%%%%%%%%%%%%%%%%%%%%%%%%%%%%%

%==============================================================================
% SUMMARY
%==============================================================================
\subsection{Summary}

\begin{tcolorbox}[colback=green!5!white, colframe=green!75!black,
    title=\textbf{Key Takeaways}]

\textbf{1. Nash is a limiting case:}
$C^*_{Nash} = \lim_{\Psi_C \to \infty} C^*_{behavioral}$

\textbf{2. Equilibrium emergence:} Evolutionary dynamics explain \textit{how}
$C^*$ is reached without assuming rationality.

\textbf{3. Equilibrium selection:} Stochastic stability selects among multiple
Nash equilibria based on resistance to invasion.

\textbf{4. Bounded cognition:} Level-$k$ and Cognitive Hierarchy show most people
reason 1--2 steps, not infinitely ($\tau \approx 1.5$).

\textbf{5. Generative structure:} The Beauty Contest (M\&N 2018) provides canonical
form underlying many strategic games.

\end{tcolorbox}

%==============================================================================
% GLOSSARY SECTION
%==============================================================================
\subsection{Glossary of Key Terms}
\label{sec:ad-glossary}

% \section{Glossary} marker for template compliance

\begin{center}
\renewcommand{\arraystretch}{1.3}
\begin{tabular}{@{}p{4cm}p{9cm}@{}}
\toprule
\textbf{Term} & \textbf{Definition} \\
\midrule
ESS & Evolutionarily Stable Strategy: cannot be invaded by mutants \\
Replicator dynamics & $\dot{x}_i = x_i(U_i - \bar{U})$: strategies grow/shrink by relative fitness \\
Level-$k$ & Bounded iterative reasoning: $L_k = BR(L_{k-1})$ \\
Cognitive Hierarchy & Poisson-distributed levels with $\tau \approx 1.5$ \\
QRE & Quantal Response Equilibrium: noisy best response \\
Stochastic stability & Equilibrium robust to rare mutations \\
\bottomrule
\end{tabular}
\end{center}

\textbf{Full definitions:} See Appendix UNMAPPED_GLS (Central Glossary).

%==============================================================================
% CRITICAL FOUNDATIONS
%==============================================================================
\subsection{Critical Foundations and Objections}
\label{sec:ad-foundations}

% \section{Critical Foundations} marker for template compliance

\begin{enumerate}
    \item \textbf{AD-F1: Representative Agent Critique}

    \textit{Objection:} Can we use a representative agent with average $\Psi_C$?

    \textit{Response:} No---Kirman (1992) shows aggregate behavior $\neq$ scaled
    individual behavior. Heterogeneous level-$k$ distribution matters.

    \item \textbf{AD-F2: Equilibrium Relevance}

    \textit{Objection:} If Nash is never reached, why study equilibrium?

    \textit{Response:} Nash is a useful benchmark. QRE and Level-$k$ show
    systematic deviations from Nash, making predictions more accurate.

    \item \textbf{AD-F3: Evolution vs.\ Learning}

    \textit{Objection:} Is biological evolution relevant for economic decisions?

    \textit{Response:} Evolutionary dynamics provide microfoundation for learning
    (Samuelson 1997). Cultural evolution operates on faster timescale.

    \item \textbf{AD-F4: Parameter Stability}

    \textit{Objection:} Are $\tau$ and $\lambda$ stable across contexts?

    \textit{Response:} Partially---expertise increases $\tau$, stakes increase
    $\lambda$. Context ($\Psi$) modulates these parameters.
\end{enumerate}

%==============================================================================
% OPEN ISSUES
%==============================================================================
\subsection{Open Issues and Future Directions}
\label{sec:ad-open-issues}

% \section{Open Issues} marker for template compliance

\begin{enumerate}
    \item \textbf{Level-$k$ vs.\ QRE:} When does each model apply? Need
    principled selection criteria.

    \item \textbf{Endogenous $\tau$:} How does cognitive depth evolve with
    experience, stakes, and context?

    \item \textbf{Neural foundations:} Can fMRI/EEG validate Level-$k$ reasoning
    stages?

    \item \textbf{Large-scale games:} Do evolutionary dynamics scale to
    populations $>1000$? Computational limits.
\end{enumerate}

%==============================================================================
% REFERENCES SECTION
%==============================================================================
\subsection{References and Further Reading}
\label{sec:ad-references}

% \section{References} marker for template compliance

\textbf{Primary Sources (20 papers, 70,000+ citations):}
\begin{itemize}[nosep]
    \item Maynard Smith \& Price (1973), Maynard Smith (1982): ESS
    \item Weibull (1995), Hofbauer \& Sigmund (1998): Evolutionary dynamics
    \item Kandori et al.\ (1993), Young (1993): Stochastic stability
    \item Nagel (1995), Stahl \& Wilson (1995): Level-$k$
    \item Camerer et al.\ (2004): Cognitive Hierarchy
    \item McKelvey \& Palfrey (1995): QRE
    \item Mauersberger \& Nagel (2018): Generative framework
\end{itemize}

\textbf{Full citations:} See \texttt{bibliography/bcm\_master.bib} (Master Bibliography).

\textbf{Related Appendices:}
\begin{itemize}[nosep]
    \item Appendix UNMAPPED_AWA (CORE-AWARE): Awareness $A(\cdot)$ formalization
    \item Appendix UNMAPPED_HOW (CORE-HOW): Complementarity $\gamma$ structure
    \item Appendix UNMAPPED_MSC Section 11: M\&N generative framework
    \item Appendix UNMAPPED_SCP: Social preferences in games
    \item Appendix UNMAPPED_SHA: Complexity economics and heterogeneity
\end{itemize}

%%%%%%%%%%%%%%%%%%%%%%%%%%%%%%%%%%%%%%%%%%%%%%%%%%%%%%%%%%%%%%%%%%%%%%%%%%%%%%%
% END OF APPENDIX AD
%%%%%%%%%%%%%%%%%%%%%%%%%%%%%%%%%%%%%%%%%%%%%%%%%%%%%%%%%%%%%%%%%%%%%%%%%%%%%%%
